\subsection{Odchylka dvou přímek v prostoru}
\label{subsec:Angle_between_lines_space}

Při zkoumání vzájemné polohy přímek jsme určovali, zda mají stejný
směr a zda mají společný bod. Nyní nás nebude zajímat jejich poloha,
ale pouze to, jaký úhel jejich směry svírají.

V rovině jsme odchylku dvou různoběžných přímek chápali jako velikost
ostrého, případně pravého úhlu, který přímky svírají. U rovnoběžných
přímek byla jejich odchylka rovna nule.

V prostoru však nastává nová situace. Dvě přímky mohou být
\textbf{mimoběžné}, a tedy nemají žádný společný bod. Nemůžeme proto
jejich odchylku určit přímo jako úhel vznikající v jejich průsečíku.

Přesto mají obě přímky jednoznačně určené směry. Odchylku mimoběžných
přímek proto můžeme určit tak, že jednu z nich rovnoběžně posuneme tak,
aby se s druhou přímkou protínala. Její směr se tím nezmění, a nezmění
se tedy ani hledaná odchylka.

\begin{figure}[H]
    \centering
    \begin{tikzpicture}[scale=0.9]



        % Bod na přímce p
        \coordinate (A) at (0,0,0);
        \coordinate (B) at (0,3,0);

        % Přímka p
        \draw[thick]
            (-3,0,0) -- (3,0,0)
            node[right] {$p$};

        % Mimoběžná přímka q
        \draw[thick]
            (-3,3,3) -- (3,3,-3)
            node[right] {$q$};

        % Rovnoběžná kopie q vedená bodem A
        \draw[dashed, thick]
            (-3,0,3) -- (3,0,-3)
            node[right] {$q'$};

        % Posunutí
        \draw[dashed] (A) -- (B);

        % Bod A
        \fill (A) circle (2pt) node[below left] {$A$};
        \fill (B) circle (2pt) node[above left] {$B$};

        % Směrové vektory
        \draw[vector]
            (A) -- (2.5,0,0)
            node[midway, below] {$\vec{u}$};

        \draw[vector]
            (A) -- (2,0,-2)
            node[midway, above right] {$\vec{v}$};

        % Symbol odchylky
        \node at (0.7,0,-0.27) {$\varphi$};

        % Úhel mezi přímkami
        \draw[-, thick] (1.25,0,0) arc (0:20:1);

        \node at (-1.1,1.1,0.2) {$q'\parallel q$};

    \end{tikzpicture}
    \caption{Odchylku mimoběžných přímek určíme pomocí rovnoběžného
    posunutí jedné z nich.}
    \label{fig:Angle_skew_lines_space}
\end{figure}

\begin{definition}
    \label{def:Angle_between_lines_space}
    \textbf{Odchylkou dvou přímek} $p$ a $q$ v prostoru nazýváme
    velikost ostrého nebo pravého úhlu mezi jejich směry.

    Odchylku označujeme $\varphi$ a platí
    \[
        \varphi\in\left\langle 0^\circ;90^\circ\right\rangle.
    \]

    Jsou-li přímky rovnoběžné nebo totožné, je jejich odchylka
    \[
        \varphi=0^\circ.
    \]
\end{definition}

Odchylka přímek tedy nezávisí na tom, kde se přímky v prostoru
nacházejí, ale pouze na jejich směrech. K jejímu výpočtu proto
potřebujeme pouze jejich směrové vektory.

Uvažujme přímky
\[
    p:\ X=A+t\vec{u}
\]
a
\[
    q:\ X=B+s\vec{v},
\]
kde
\[
    \vec{u}=(u_x;u_y;u_z)
    \qquad\text{a}\qquad
    \vec{v}=(v_x;v_y;v_z)
\]
jsou jejich nenulové směrové vektory.

Z kapitoly o skalárním součinu vektorů
(viz vztah \ref{eq:Angle_between_vectors_space}) víme, že pro úhel
$\alpha$ mezi dvěma nenulovými vektory platí
\[
    \cos\alpha
    =
    \frac{\vec{u}\cdot\vec{v}}
         {|\vec{u}|\cdot|\vec{v}|}.
\]

Úhel dvou vektorů však může nabývat hodnot od $0^\circ$ do
$180^\circ$, zatímco odchylka přímek je vždy nejvýše $90^\circ$.

Důvodem je skutečnost, že přímka nemá žádnou pevně určenou orientaci.
Je-li $\vec{u}$ směrovým vektorem přímky, potom je jejím směrovým
vektorem také opačný vektor
\[
    -\vec{u}.
\]

Pokud například směrové vektory svírají úhel
\[
    \alpha=120^\circ,
\]
pak po obrácení jednoho z nich získáme úhel
\[
    180^\circ-120^\circ=60^\circ.
\]
Právě menší z těchto dvou úhlů představuje odchylku přímek.

Platí tedy
\[
    \varphi
    =
    \min\{\alpha;180^\circ-\alpha\}.
\]

Protože
\[
    \cos(180^\circ-\alpha)=-\cos\alpha,
\]
můžeme oba případy zahrnout do jediného vztahu pomocí absolutní
hodnoty:

\begin{equation}
        \cos\varphi
        =
        \frac{
            \left|\vec{u}\cdot\vec{v}\right|
        }{
            |\vec{u}|\cdot|\vec{v}|
        }
    \label{eq:Angle_between_lines_space}
\end{equation}

Po rozepsání do souřadnic dostáváme

\begin{equation}
    \cos\varphi
    =
    \frac{
        \left|
            u_xv_x+u_yv_y+u_zv_z
        \right|
    }{
        \sqrt{u_x^2+u_y^2+u_z^2}
        \sqrt{v_x^2+v_y^2+v_z^2}
    }
    \label{eq:Angle_between_lines_space_coordinates}
\end{equation}

a samotnou odchylku vypočítáme jako
\[
    \varphi
    =
    \arccos
    \left(
        \frac{
            \left|\vec{u}\cdot\vec{v}\right|
        }{
            |\vec{u}|\cdot|\vec{v}|
        }
    \right).
\]

\begin{remark}
    \label{rem:Absolute_value_angle_lines_space}
    Absolutní hodnota ve vztahu
    \[
        \cos\varphi
        =
        \frac{|\vec{u}\cdot\vec{v}|}
             {|\vec{u}|\cdot|\vec{v}|}
    \]
    je důležitá.

    Nahradíme-li směrový vektor $\vec{u}$ opačným vektorem
    $-\vec{u}$, přímka se nijak nezmění, ale skalární součin změní
    znaménko:
    \[
        (-\vec{u})\cdot\vec{v}
        =
        -(\vec{u}\cdot\vec{v}).
    \]

    Odchylka přímek však na volbě orientace jejich směrových vektorů
    záviset nesmí. Absolutní hodnota tuto závislost odstraní a současně
    zajistí, že
    \[
        \cos\varphi\geq0,
    \]
    takže výsledná odchylka leží v intervalu
    \[
        \left\langle0^\circ;90^\circ\right\rangle.
    \]
\end{remark}

Při výpočtu odchylky tedy není nutné hledat průsečík přímek ani
předem určovat jejich vzájemnou polohu. Z parametrických vyjádření
stačí určit jejich směrové vektory.

\begin{example}
    \label{exa:Angle_between_intersecting_lines_space}

    Určete odchylku přímek
    \[
        \begin{aligned}
            p: x&=1+t,\\
            y&=-1+t,\\
            z&=2; t\in\mathbb{R}
        \end{aligned}
    \]
    a
    \[
        \begin{aligned}
            q: x&=3-s,\\
            y&=2,\\
            z&=1-s; s\in\mathbb{R}.
        \end{aligned}
    \]

    Směrové vektory přímek jsou
    \[
        \vec{u}=(1;1;0),
        \qquad
        \vec{v}=(-1;0;-1).
    \]

    Nejprve vypočítáme jejich skalární součin:
    \[
        \vec{u}\cdot\vec{v}
        =
        1\cdot(-1)+1\cdot0+0\cdot(-1)
        =
        -1.
    \]

    Velikosti směrových vektorů jsou
    \[
        |\vec{u}|
        =
        \sqrt{1^2+1^2+0^2}
        =
        \sqrt{2}
    \]
    a
    \[
        |\vec{v}|
        =
        \sqrt{(-1)^2+0^2+(-1)^2}
        =
        \sqrt{2}.
    \]

    Dosadíme do vztahu pro odchylku přímek:
    \begin{align*}
        \cos\varphi
        &=
        \frac{|\vec{u}\cdot\vec{v}|}
             {|\vec{u}|\cdot|\vec{v}|}\\
        &=
        \frac{|-1|}
             {\sqrt{2}\cdot\sqrt{2}}\\
        &=
        \frac12.
    \end{align*}

    Odtud
    \[
        \varphi=60^\circ.
    \]

    Kdybychom absolutní hodnotu nepoužili, získali bychom
    \[
        \cos\alpha=-\frac12
    \]
    a tedy úhel směrových vektorů
    \[
        \alpha=120^\circ.
    \]

    Odchylkou přímek je však menší z obou úhlů:
    \[
        180^\circ-120^\circ=60^\circ.
    \]
\end{example}

\begin{remark}
    \label{rem:Special_angles_lines_space}
    Ze vztahu pro odchylku můžeme okamžitě určit dva důležité
    speciální případy.

    Jsou-li směrové vektory rovnoběžné, tedy
    \[
        \vec{v}=k\vec{u},
        \qquad k\neq0,
    \]
    potom
    \[
        \varphi=0^\circ.
    \]

    To nastává jak u rovnoběžných, tak u totožných přímek.

    Je-li naopak
    \[
        \vec{u}\cdot\vec{v}=0,
    \]
    potom
    \[
        \cos\varphi=0
    \]
    a tedy
    \[
        \varphi=90^\circ.
    \]

    Směry přímek jsou v tomto případě navzájem kolmé.

    Samotná hodnota odchylky však neurčuje vzájemnou polohu přímek.
    Dvě přímky se stejnou nenulovou odchylkou mohou být různoběžné,
    nebo mimoběžné.
\end{remark}

Právě poslední skutečnost je jedním z důležitých rozdílů oproti
geometrii v rovině. V rovině se dvě přímky s různými směry vždy
protínají. V prostoru tomu tak být nemusí.

\begin{example}
    \label{exa:Angle_skew_lines_ninety_degrees}

    Uvažujme přímky
    \[
        \begin{aligned}
            p: x&=2+t,\\
            y&=3-2t,\\
            z&=-1+t; t\in\mathbb{R}
        \end{aligned}
    \]
    a
    \[
        \begin{aligned}
            q: x&=-2+3s,\\
            y&=3+2s,\\
            z&=3+s; s\in\mathbb{R}.
        \end{aligned}
    \]

    V příkladu \ref{exa:Skew_parametric_lines_space} jsme zjistili,
    že tyto přímky jsou mimoběžné.

    Jejich směrové vektory jsou
    \[
        \vec{u}=(1;-2;1),
        \qquad
        \vec{v}=(3;2;1).
    \]

    Vypočítáme jejich skalární součin:
    \begin{align*}
        \vec{u}\cdot\vec{v}
        &=
        1\cdot3+(-2)\cdot2+1\cdot1\\
        &=
        3-4+1\\
        &=
        0.
    \end{align*}

    Proto
    \[
        \cos\varphi=0
    \]
    a tedy
    \[
        \varphi=90^\circ.
    \]

    Přímky jsou mimoběžné a přesto mají navzájem kolmé směry.
    To znamená, že odchylka dvou mimoběžných přímek může být také
    pravý úhel.
\end{example}

\begin{remark}
    \label{rem:Perpendicular_directions_skew_lines}
    Je třeba rozlišovat kolmost směrů a existenci průsečíku.

    Podmínka
    \[
        \vec{u}\cdot\vec{v}=0
    \]
    říká pouze to, že odchylka přímek je $90^\circ$. Neříká nic o tom,
    zda se přímky skutečně protínají.

    Pokud se přímky navíc protínají, můžeme je chápat jako dvě
    klasické kolmé různoběžky. Pokud společný bod nemají, mohou být
    mimoběžné s odchylkou $90^\circ$.
\end{remark}

Směrové vektory nemusíme vždy získávat přímo z parametrického
vyjádření. Je-li přímka určena dvěma body, použijeme jako její směrový
vektor rozdíl jejich souřadnic.

\begin{example}
    \label{exa:Angle_lines_defined_by_points_space}

    Jsou dány body $A[1;0;2], B[2;1;2],C[-1;2;0],D[0;2;1]$.
    Určete odchylku přímek $AB$ a $CD$.

    Směrový vektor přímky $AB$ je
    \[
        \overrightarrow{AB}
        =
        B-A
        =
        (1;1;0).
    \]

    Směrový vektor přímky $CD$ je
    \[
        \overrightarrow{CD}
        =
        D-C
        =
        (1;0;1).
    \]

    Jejich skalární součin je
    \[
        \overrightarrow{AB}\cdot\overrightarrow{CD}
        =
        1\cdot1+1\cdot0+0\cdot1
        =
        1.
    \]

    Dále
    \[
        |\overrightarrow{AB}|=\sqrt2,
        \qquad
        |\overrightarrow{CD}|=\sqrt2.
    \]

    Proto
    \[
        \cos\varphi
        =
        \frac{1}{\sqrt2\cdot\sqrt2}
        =
        \frac12.
    \]

    Odchylka přímek je
    \[
        \varphi=60^\circ.
    \]
\end{example}

V některých úlohách nemusí být odchylka neznámou veličinou. Naopak
může být zadána a naším úkolem je určit neznámý parametr tak, aby
směrové vektory danou odchylku skutečně svíraly.

\begin{example}
    \label{exa:Angle_lines_parameter_space}

    Je dána přímka
    \[
        \begin{aligned}
            p: x&=1+t,\\
            y&=2+t,\\
            z&=3; t\in\mathbb{R}
        \end{aligned}
    \]
    a soustava přímek
    \[
        \begin{aligned}
            q_m: x&=-1+s,\\
            y&=4+ms,\\
            z&=2+s; s\in\mathbb{R}.
        \end{aligned}
    \]

    Určete hodnotu parametru $m\in\mathbb{R}$ tak, aby odchylka
    přímek $p$ a $q_m$ byla $45^\circ$.

    Směrový vektor přímky $p$ je
    \[
        \vec{u}=(1;1;0),
    \]
    zatímco směrový vektor přímky $q_m$ je
    \[
        \vec{v}=(1;m;1).
    \]

    Vypočítáme potřebné hodnoty:
    \[
        \vec{u}\cdot\vec{v}
        =
        1+m,
    \]
    \[
        |\vec{u}|=\sqrt2
    \]
    a
    \[
        |\vec{v}|
        =
        \sqrt{1+m^2+1}
        =
        \sqrt{m^2+2}.
    \]

    Protože požadujeme odchylku $45^\circ$, musí platit
    \[
        \frac{|1+m|}
             {\sqrt2\sqrt{m^2+2}}
        =
        \cos45^\circ
        =
        \frac{\sqrt2}{2}.
    \]

    Po vynásobení obou stran dostáváme
    \[
        |1+m|
        =
        \sqrt{m^2+2}.
    \]

    Obě strany jsou nezáporné, můžeme proto rovnici umocnit:
    \begin{align*}
        (1+m)^2&=m^2+2,\\
        1+2m+m^2&=m^2+2,\\
        2m&=1,\\
        m&=\frac12.
    \end{align*}

    Ověříme získanou hodnotu:
    \[
        \vec{v}
        =
        \left(1;\frac12;1\right).
    \]

    Potom
    \[
        \vec{u}\cdot\vec{v}
        =
        1+\frac12
        =
        \frac32
    \]
    a
    \[
        |\vec{v}|
        =
        \sqrt{1+\frac14+1}
        =
        \frac32.
    \]

    Proto
    \[
        \cos\varphi
        =
        \frac{\frac32}
             {\sqrt2\cdot\frac32}
        =
        \frac{1}{\sqrt2}
        =
        \frac{\sqrt2}{2}.
    \]

    Hledaná hodnota parametru je tedy $m=\frac{1}{2}$.
    Přímka $q$ má poté parametrické vyjádření:
    \begin{align*}
        q:  x&=-1+s,\\
            y&=4+\frac12s,\\
            z&=2+s; s\in\mathbb{R}.
    \end{align*}

\end{example}

\begin{remark}
    \label{rem:Angle_line_coordinate_axes_space}
    Za přímky můžeme považovat také jednotlivé souřadnicové osy.
    Jejich směrové vektory můžeme zvolit
    \[
        \vec{e}_x=(1;0;0),
        \qquad
        \vec{e}_y=(0;1;0),
        \qquad
        \vec{e}_z=(0;0;1).
    \]

    Je-li
    \[
        \vec{u}=(u_x;u_y;u_z)
    \]
    směrový vektor přímky $p$, potom pro její odchylky
    $\alpha$, $\beta$ a $\gamma$ od os $x$, $y$ a $z$ platí
    \[
        \cos\alpha
        =
        \frac{|u_x|}{|\vec{u}|},
    \]
    \[
        \cos\beta
        =
        \frac{|u_y|}{|\vec{u}|},
    \]
    \[
        \cos\gamma
        =
        \frac{|u_z|}{|\vec{u}|}.
    \]

    Z toho zároveň vyplývá
    \[
        \cos^2\alpha
        +
        \cos^2\beta
        +
        \cos^2\gamma
        =
        1.
    \]
\end{remark}

\begin{remark}
    \label{rem:Infinitely_many_lines_given_angle_space}
    V prostoru se oproti rovině objevuje ještě jeden důležitý rozdíl.

    V rovině můžeme daným bodem vést zpravidla pouze dvě přímky,
    které svírají s danou přímkou zadanou nenulovou odchylku
    $\varphi$.

    V prostoru však existuje takových přímek nekonečně mnoho.
    Jejich směrové vektory mohou kolem směru původní přímky měnit
    svou polohu, přičemž stále zachovávají stejnou odchylku.

    Pro
    \[
        0^\circ<\varphi<90^\circ
    \]
    vytvářejí všechny tyto směry povrch kužele, jehož osou je směr
    dané přímky.

    Pro
    \[
        \varphi=90^\circ
    \]
    získáme nekonečně mnoho kolmých směrů. Přímky vedené jedním bodem
    v těchto směrech společně vyplňují rovinu kolmou k původní přímce.
    S rovinami a jejich analytickým popisem se budeme zabývat později.
\end{remark}

Celý postup při určování odchylky dvou přímek v prostoru můžeme
shrnout následovně. Z parametrických vyjádření nebo ze zadaných bodů
nejprve určíme nenulové směrové vektory
\[
    \vec{u}
    \qquad\text{a}\qquad
    \vec{v}.
\]
Poté vypočítáme jejich skalární součin a velikosti a použijeme vztah
\[
    \cos\varphi
    =
    \frac{
        |\vec{u}\cdot\vec{v}|
    }{
        |\vec{u}|\cdot|\vec{v}|
    }.
\]

Odchylka přímek závisí pouze na jejich směrech. Proto není při jejím
výpočtu nutné znát jejich průsečík ani rozhodovat, zda jsou přímky
různoběžné, rovnoběžné, totožné nebo mimoběžné.